\documentclass[11pt,a4paper]{article}
\usepackage[margin=1in]{geometry}
\usepackage{hyperref}
\usepackage{enumitem}
\usepackage{graphicx}
\usepackage{xcolor}

% -----------------------------------------------------
% bvraghav-tiet-ucs503p-article.sty
% -----------------------------------------------------

% Variable: Lab Instructor
\makeatletter
\newcommand{\BvrSetLabInstructor}[1]{%
  \gdef\Bvr@LabInstructor{#1}}
\newcommand{\Bvr@LabInstructor}{Dr. Raghav B. Venkataramaiyer}
\newcommand{\BvrLabInstructorValue}{\Bvr@LabInstructor}
\makeatother

% Add subtitle
\usepackage{titling}

% Set title bold
\pretitle{\begin{center}\bfseries\fontsize{18bp}{18bp}\selectfont}

\makeatletter
\newcommand{\BvrSubtitle}[1]{%
  \posttitle{%
    \par\end{center}
    \begin{center}\large#1\end{center}
    \vskip 0.5em}%
}
\makeatother

% Hints in the section title(s)
\newcommand{\BvrHint}[1]{%
  {\normalfont\normalsize\itshape\hfill #1}}

% -----------------------------------------------------

\title{MeterIQ: Explainable Smart Meter Analytics Platform}

\BvrSetLabInstructor{Dr. Jeelani}

\BvrSubtitle{%
  Project Proposal for UCS503P\\
  Submitted to: \BvrLabInstructorValue\\
  \bfseries Thapar Institute of Engineering and Technology%
}

\author{
  Yashit Arora, CSED%
  \thanks{Roll No: \texttt{1024160006}, Email: \texttt{yarora1\_be24@thapar.edu}}
  \and
  Aarushi Gahlawat, CSED%
  \thanks{Roll No: \texttt{1024160008}, Email: \texttt{agahlawat\_be24@thapar.edu}}
}

\date{\today}

\begin{document}

\maketitle

\section{Higher-order goal%
  \BvrHint{\ldots data-driven electricity monitoring}}

MeterIQ aims to make smart-meter data easier to analyse and use.
The system will identify unusual consumption behaviour and provide
short-term load forecasts.

We will work with smart-meter datasets from Mathura and Bareilly.
The main goal is to bring anomaly detection and load forecasting
together in one platform so that a user can view a meter's
consumption, investigate unusual behaviour and see its expected
future load.

\section{Time-to-value %
  \BvrHint{\ldots primary heuristic}}

\begin{itemize}[leftmargin=0.7cm]

    \item Understand and clean the Mathura and Bareilly datasets.

    \item Build simple baselines for anomaly detection and load
    forecasting.

    \item Compare suitable ML approaches using the prepared data.

    \item Connect the analysis with a backend, data storage and
    dashboard.

    \item Use Git and GitHub to track development.

\end{itemize}

\section{Problem Statement}

Smart meters generate a large amount of electricity data, but
manually checking this data for unusual behaviour is difficult.
Sudden changes in consumption may have different causes, and it is
not practical to inspect every reading individually.

The available datasets contain timestamped consumption along with
measurements such as voltage, current and frequency. The project
aims to use these readings to identify unusual patterns and forecast
future load in a way that is easier for a user to understand.

An anomaly will be treated as a signal for further investigation,
not as automatic proof of electricity theft.

\section{Proposed Solution %
  \BvrHint{\ldots Software Proposition}}

\subsection{Overview}

MeterIQ will be developed as a web-based smart-meter analytics
platform with two main functions: anomaly detection and short-term
load forecasting.

The anomaly module will identify unusual meter behaviour and provide
an anomaly or risk indication. The forecasting module will estimate
future electricity consumption from historical readings.

Both results will be available through the same application instead
of being handled as separate analysis notebooks.

\subsection{Core workflow}

\begin{enumerate}[leftmargin=0.7cm]

    \item Load and inspect the smart-meter data.

    \item Clean the data and prepare useful features.

    \item Run anomaly detection and load forecasting.

    \item Store and expose the results through the application.

    \item Allow the user to search for a meter and view its history,
    alerts and forecast.

\end{enumerate}

\subsection{Operational constraints for an educational Software Engineering project}

\begin{itemize}[leftmargin=0.7cm]

    \item The actual dataset will be studied before selecting
    features and models.

    \item Chronological splits will be used for forecasting to avoid
    data leakage.

    \item Anomaly results will be treated as investigation
    indicators.

    \item The system will be developed in separate modules so that
    the parts can be tested independently.

\end{itemize}

\section{Solution Approach %
  \BvrHint{\ldots Engineering focus}}

The project will follow a simple workflow: understand the data,
prepare it, build the ML components and then connect them to the
application.

\subsection{Data understanding and preprocessing %
  \BvrHint{\ldots data quality}}

The Mathura and Bareilly datasets will be inspected separately. The
available data includes timestamp, energy consumption, average
voltage, average current, frequency and meter information.

Preprocessing will include timestamp handling, checking missing and
duplicate records, checking unusual values and validating meter-wise
data. Time-based, lag and rolling features will be created where
useful.

\subsection{Anomaly and risk detection %
  \BvrHint{\ldots ML component}}

The anomaly module will look for consumption behaviour that differs
from a meter's usual pattern. Voltage, current, frequency and
time-based information may also be considered.

Models such as Isolation Forest and Local Outlier Factor will be
compared. The final approach will be selected after testing it on
the available data.

The output will include an anomaly indication and supporting
information so that the user can understand why a reading was
flagged.

\subsection{Short-term load forecasting %
  \BvrHint{\ldots forecasting component}}

The forecasting module will use historical consumption to estimate
future load. A simple historical baseline will be used first and
then compared with ML models such as Linear Regression and Random
Forest.

Lagged consumption, rolling statistics and calendar features will be
considered. Models will be evaluated using MAE, RMSE and MAPE on a
chronological test period.

\subsection{Explainability %
  \BvrHint{\ldots understandable results}}

MeterIQ will show more than an anomaly flag. For a flagged reading,
the user will be able to see the consumption pattern and relevant
changes that contributed to the alert.

The system will clearly separate model output from the information
calculated from the data so that an anomaly is not presented as a
confirmed cause.

\subsection{Web architecture %
  \BvrHint{\ldots web-based solution}}

The application will have a web-based frontend, a backend layer,
data-processing and ML components, and a suitable data-storage
solution.

Python, Pandas, NumPy and Scikit-learn will be used for the data and
ML work. Git and GitHub will be used for version control.

\subsection{CI/CD and fast delivery enabler}

GitHub will be used throughout development. Automated testing and
basic build checks can be added as the application develops to make
changes easier to test and manage.

\section{Evaluation Criterion %
  \BvrHint{\ldots measurable and attributable}}

\subsection{Primary evaluation metric}

For load forecasting, the main evaluation metrics will be:

\begin{itemize}[leftmargin=0.7cm]

    \item Mean Absolute Error (MAE)

    \item Root Mean Squared Error (RMSE)

    \item Mean Absolute Percentage Error (MAPE)

\end{itemize}

For anomaly detection, the detected patterns will be checked for
consistency and usefulness. Confirmed theft labels will not be
assumed unless they are available in the data.

\subsection{Secondary evaluation metrics}

The application will also be checked for correct data retrieval,
working API calls, correct display of results, basic response time
and integration between the ML and application components.

\subsection{Pilot validation plan %
  \BvrHint{\ldots high-level}}

Forecasting models will be tested chronologically so that future
information is not used during training. Different models will be
compared using the same evaluation period.

Anomaly results will be inspected at the meter level and compared
with historical behaviour.

\section{Scalability %
  \BvrHint{\ldots with foresight}}

The system will be designed so that additional meters and future
readings can be added later. Keeping the data, ML and application
parts separate will also make it easier to update the models when
more data becomes available.

\section{Engine availability heuristic}

The availability of a meter will be checked using the completeness
and continuity of its readings. Large gaps in the data may affect
forecasting and anomaly analysis, so data availability will be
considered before using a meter for detailed analysis.

\section{Project scope and deliverables %
  \BvrHint{\ldots iteration-friendly}}

\subsection{Initial deliverable %
  \BvrHint{\ldots first iteration}}

The first deliverable will be a cleaned and documented version of
the Mathura and Bareilly data with exploratory analysis of
consumption and electrical behaviour.

\subsection{Subsequent deliverables}

\begin{itemize}[leftmargin=0.7cm]

    \item Data preprocessing and feature engineering pipeline.

    \item Anomaly detection module.

    \item Short-term load forecasting module.

    \item Model evaluation and comparison.

    \item Explainability for anomaly results.

    \item Backend and data-storage integration.

    \item Web dashboard for meter analysis.

    \item Testing and project documentation.

\end{itemize}

\section{Risks and mitigations}

\begin{itemize}[leftmargin=0.7cm]

    \item \textbf{Data quality:} missing, duplicate or unusual
    readings will be checked during preprocessing.

    \item \textbf{Different meter behaviour:} analysis will consider
    meter-level patterns instead of assuming all meters behave the
    same way.

    \item \textbf{False anomalies:} flagged readings will be treated
    as investigation signals rather than confirmed theft.

    \item \textbf{Data leakage:} chronological validation and
    carefully created lag features will be used.

    \item \textbf{Model selection:} simple baselines will be compared
    with ML models before choosing the final approach.

    \item \textbf{System integration:} the ML, backend and frontend
    components will be developed separately and integrated step by
    step.

\end{itemize}

\section{Summary}

MeterIQ will analyse smart-meter data from Mathura and Bareilly using
two connected components: anomaly detection and short-term load
forecasting.

The system will help a user view meter consumption, identify unusual
behaviour and understand the information behind an alert. It will
also provide future load estimates based on historical readings.

The final project will combine the data pipeline, ML models and
web-based application into one system. The focus will be on useful
and understandable results rather than making unsupported claims
about electricity theft.

\end{document}
